\chapter{STL、区间和视图}

STL是C++标准库的一部分，也可以认为是最重要的一部分。它提供了许多通用的数据结构和算法，极大地提高了C++程序的开发效率。

区间和视图是C++标准库最近新增的一部分，它们提供了一种更高层次的抽象，使得操作集合和序列变得更加简洁和直观。区间（Ranges）允许我们以一种声明式的方式处理数据，而视图（Views）则提供了一种惰性求值的机制，使得我们可以在不创建中间容器的情况下对数据进行操作。如果C++20成为默认的处理标准，那么区间和视图将成为C++程序设计中不可或缺的工具。

本章C++标准为C++23。

\section{C++的STL}

STL分三部分：容器（Containers）、算法（Algorithms）和迭代器（Iterators）。容器是用来存储数据的对象，算法是对数据进行操作的函数，而迭代器则是连接容器和算法的桥梁。通过迭代器，算法可以在不同类型的容器上工作，而不需要了解容器的具体实现。但要写好C++，还是要了解一下容器的实现原理，以及相应算法的复杂度。因此，本章会以容器为纲目，展开介绍算法和迭代器，并在介绍的过程中，穿插一些算法和数据结构的基础知识。

STL的容器分为三类：相同类型元素的容器、不同类型元素的容器、其他容器。

相同类型元素的容器最多，它们可以存储复数个相同类型的元素。
\begin{itemize}
    \item 顺序表：\verb|std::array|、\verb|std::vector|、\verb|std::deque|、\verb|std::list|、\verb|std::forward_list|；
    \item 关联容器：\verb|std::set|、\verb|std::map|、\verb|std::multiset|、\verb|std::multimap|；
    \item 无序关联容器：\verb|std::unordered_set|、\verb|std::unordered_map|、\verb|std::unordered_multiset|、\verb|std::unordered_multimap|；
    \item 容器适配器：\verb|std::stack|、\verb|std::queue|、\verb|std::priority_queue|；
    \item 字符串：\verb|std::string|等；
    \item 位集合：\verb|std::bitset|。
\end{itemize}

不同类型元素的容器较少，它们可以存储复数个不同类型的元素。只有两个：对 \verb|pair|，元组 \verb|tuple|。

最后是其他容器，主要是一些特殊用途的容器，如变体 \verb|variant|，可选值 \verb|optional|，任意类型 \verb|any|，智能指针 \verb|shared_ptr|、\verb|unique_ptr|、\verb|weak_ptr|。

不同容器的不同使用情况见下表。

\begin{longtable}[c]{l|l|p{7cm}}
    \caption{STL容器底层与使用场景简表}\label{tab:stl-containers}\\
    \toprule
    底层 & 容器 & 使用场景 \\
    \midrule
    \endfirsthead          % 首页表头

    \multicolumn{3}{c}{\footnotesize 续表~ \ref{tab:stl-containers} }\\[.5ex]
    \toprule
    底层 & 容器 & 使用场景 \\
    \midrule
    \endhead               % 后续页表头

    \midrule
    \multicolumn{3}{r}{\footnotesize 接下页}
    \endfoot               % 每页底部（除末页）

    \bottomrule
    \endlastfoot           % 末页底部

    \multirow{2}{*}{顺序表} & \verb|std::array| & 数据全在栈上，性能较为敏感，适用于编译期的定长小缓存：存储固定的配置表、常量数据、小矩阵等。 \\
    & \verb|std::vector| & 通用首选，各种场景都可以使用，尤其是频繁遍历、尾部增删的场景 \\
    \midrule
    \multirow{3}{*}{分块顺序表} & \verb|std::deque| & 适用于两头进出但中间不常访问的场景，例如消息队列、撤销重做历史记录等 \\
    & \verb|std::stack| & LIFO逻辑，主要用于模拟递归等。 \\
    & \verb|std::queue| & FIFO逻辑，主要用于模拟缓存等。 \\
    \midrule
    带SSO的顺序表 & \verb|std::string| & 一切文本处理。小字符串会存在栈上，大字符串存在堆上。切记：为节省开销，如不修改，尽量声明为常量。 \\
    \midrule
    单链表 & \verb|std::forward_list| & 内存极度受限、中间要插入和删除、且无需反向遍历的海量节点存储器。 \\
    \midrule
    双链表 & \verb|std::list| & 频繁中间插入和删除的场景，如LRU缓存等。 \\
    \midrule
    \multirow{4}{*}{红黑树} & \verb|std::set| & 自动去重排序 \\
    & \verb|std::map| & 有序键值字典 \\
    & \verb|std::multiset| & 允许少量重复元素的有序集合，多重复元素慎用 \\
    & \verb|std::multimap| & 慎用，建议回退到 \verb|std::map| \\
    \midrule
    \multirow{4}{*}{哈希表} & \verb|std::unordered_set| & 高频查找的快速去重和查询，只关心存在性 \\
    & \verb|std::unordered_map| & 需要高频查找的键值对缓存，不关心顺序 \\
    & \verb|std::unordered_multiset| & 允许重复元素的快速计数集合 \\ 
    & \verb|std::unordered_multimap| & 存在重大设计缺陷，不要使用。建议回退到 \verb|std::unordered_map| \\
    \midrule
    堆 & \verb|std::priority_queue| & 海量数据top-k筛选 \\
    \midrule
    \multirow{2}{*}{位集合} & \verb|std::bitset| & 编译期定长标志位，硬件多用 \\
    & \verb|std::vector<bool>| & 动态大位图，多用于内存压缩。行为较异常，禁止在多线程使用，内存充足时可以考虑回退到 \verb|std::vector<unsigned char>| \\
\end{longtable}

\subsection{前置知识：迭代器和STL算法}

迭代器是STL所有容器都有的对象，它们提供了一种统一的方式来访问容器中的元素。迭代器可以看作是指向容器中元素的指针，但它们不仅仅是指针，还提供了许多额外的功能，如遍历、修改和比较等。

比如说，对于一个 \verb|std::vector<int> vec|，有
\begin{minted}{cpp}
auto it = vec.begin(); // it 是一个迭代器，指向 vec 的第一个元素
auto it2 = vec.end(); // it2 是一个迭代器，指向 vec 的最后一个元素的下一个位置
for (auto it = vec.begin(); it != vec.end(); ++it) { // 迭代器可以算术加减
    std::cout << *it << " "; // 通过解引用迭代器访问元素
}
\end{minted}

上述循环也可以写成以下两种形式：
\begin{minted}{cpp}
for(auto it : vec) { // 范围for循环，自动使用迭代器
    std::cout << it << " "; // 直接访问元素
}
std::for_each(vec.begin(), vec.end(), [](int x) { // 使用算法 std::for_each
    std::cout << x << " "; // 通过lambda函数访问元素
});
\end{minted}

迭代器确实可以算术加减，包括：
\begin{itemize}
    \item 自增、自减、普通的加减。加（减）一个整数 $n$，迭代器会向前（向后）移动 $n$ 个元素。
    \item 迭代器之间的差值。两个迭代器相减，得到的是它们之间的元素个数。因此，想知道某迭代器 \verb|it| 在容器中的索引，可以用 \verb|it - vec.begin()| 来计算。
    \item 迭代器的比较。两个迭代器可以进行大小比较，判断它们在容器中的相对位置。
\end{itemize}
迭代器的算术加减也可以用 \verb|std::advance(it, n)| 来实现，它会将迭代器 \verb|it| 向前（向后）移动 $n$ 个元素。

除了最基础的 \verb|begin()| 和 \verb|end()|，STL的容器还提供了其他一些迭代器，如 \verb|rbegin()|、\verb|rend()|、\verb|cbegin()|、\verb|cend()|、\verb|crbegin()|、\verb|crend()| 等。带字母 \verb|r| 的是反向迭代器，表示从后往前遍历；带字母 \verb|c| 的是常量迭代器，表示不能修改元素的值。

迭代器从底层实现上，分为随机访问迭代器和顺序访问迭代器。对于支持随机访问的数据类型，如顺序表，迭代器可以直接跳转到任意位置（即 \verb|&vec[x] + y * sizeof(T)|）；而对于不支持随机访问的数据类型，如链表，迭代器必须一个一个的从头遍历到尾，不能直接跳转到任意位置。

对于绝大多数算法，有两种版本：第一种是利用迭代器的普通算法，另一种是约束版本的算法（C++20）。例如要在一个表中查找第一个等于 $x$ 的元素：
\begin{minted}{cpp}
auto it = std::find(vec.begin(), vec.end(), x); // 普通算法，返回迭代器
auto it2 = std::ranges::find(vec, x); // 约束版本的算法，返回迭代器
\end{minted}

部分算法允许原生多线程，适用于非常大的数据结构。当数据结构不太大时，建议不使用原生多线程。这要求应用该算法的容器必须采用随机访问迭代器（如顺序表）。其语法就是在函数后面加一个 \verb|std::execution::xx|，其中 \verb|xx| 可以是 \verb|seq|、\verb|par|、\verb|par_unseq|，分别表示顺序执行、并行执行、并行无序执行。在数据结构大小超过50万后，建议使用 \verb|par|。虽然 \verb|par_unseq| 理论上有着最快的速度，但对于不少算法——例如 \verb|std::sort|，它的实现是基于 \verb|std::stable_sort| 的，而 \verb|std::stable_sort| 是不支持 \verb|par_unseq| 的——会退化为 \verb|par|，它只在数据结构极大、对数据之间的访问独立（如 \verb|std::transform|）时，才会有明显的速度优势。因此，除非你非常了解算法的实现，否则不建议使用 \verb|par_unseq|。在数据结构小于一万时，建议使用 \verb|seq| 或不写（默认为 \verb|seq|）。此外，Clang对这个写法的支持稍差，建议使用GCC编译之。

\subsection{顺序表}

顺序表指的是在内存中连续、线性存储的元素集合。下一个元素的地址一定是上一个元素的地址加上等于元素大小的偏移量。

这两个容器，我们在C++基础中已经简要的介绍了：
\begin{itemize}
    \item \verb|std::array<T, N>| 是固定大小为 \verb|N| 的数组，且固定在栈上分配内存。它的大小在编译时就已经确定，不能动态改变。
    \item \verb|std::vector<T>| 是动态大小的数组，大小可以随时扩容（但一般不会缩容），它的内存分配在堆上进行。它的大小可以在运行时动态改变。
\end{itemize}
接下来我们将以 \verb|std::vector<int> vec| 为例。

\verb|std::vector| 的构造方法，最常见的有四种：
\begin{minted}{cpp}
    std::vector<int> vec1; // 默认构造函数，创建一个空的 vector
    std::vector<int> vec2(10); // 创建一个大小为 10 的 vector，元素默认初始化为 0
    std::vector<int> vec3(10, 5); // 创建一个大小为 10 的 vector，并将所有元素初始化为 5
    std::vector<int> vec4 = {1, 2, 3, 4, 5}; // 直接从列表初始化
\end{minted}

顺序表最基础的几个操作的复杂度包括：
\begin{itemize}
    \item 访问指定索引的元素所用的时间复杂度为 $O(1)$。这是显然的，只需要一次计算就可以得到对应元素的地址。对应的算法是 \verb|vec.at(i)| 或者 \verb|vec[i]|。建议使用 \verb|vec.at(i)|，因为它会进行边界检查，如果索引越界，会抛出异常，而后者不会进行边界检查，可能会导致未定义行为。
    \item 在顺序表的末尾添加一个元素所用的时间复杂度为 $O(1)$（不考虑扩容）。如果添加这个元素引发了扩容，那么时间复杂度为 $O(n)$，其中 $n$ 是当前顺序表的大小。因为需要分配一块新的内存，并将原有元素复制到新内存中。对应的算法是 \verb|vec.push_back(x)| 和 \verb|vec.emplace_back(x)|：前者会将一个“已经构造好”的 \verb|x| 添加到顺序表的末尾；后者会把 \verb|x| 视为构造 \verb|T| 的参数，直接在顺序表的末尾构造一个 \verb|T| 对象。在已有对象时，两者性能差不太多；在需要构造时，后者更快，因为前者需要构造和移动两步操作。
    \item 在顺序表的末尾删除一个元素所用的时间复杂度为 $O(1)$。对应的算法是 \verb|vec.pop_back()|。
    \item 在顺序表的中间或开头插入或删除一个元素，所用的时间复杂度为$O(n)$，其中 $n$ 是当前顺序表的大小。因为需要移动后续的元素以腾出空间或填补空缺。对应的算法是 \verb|vec.insert(it, x)| 和 \verb|vec.erase(it)|，其中 \verb|it| 是一个迭代器，指向要插入或删除的位置。
\end{itemize}

什么时候用顺序表？实际上大多数情况用顺序表都是没有问题的，除非：
\begin{itemize}
    \item 频繁的在中间插入删除；
    \item 需要其他的特殊功能，如双端队列在头部的插删、优先队列最大（小）值始终在队首等。
\end{itemize}

我们将以顺序表为例，介绍最常用的一些算法。其他的容器，也有类似的算法。对于适用的算法，我们不会再在之后的章节重复介绍。

\paragraph{本身的性质} 用以下操作可以获取或更改顺序表的性质，它们都是 \verb|std::vector| 的成员函数：
\begin{itemize}
    \item \verb|vec.size()| 返回顺序表的大小，即当前存储的元素个数；
    \item \verb|vec.capacity()| 返回顺序表的容量，即当前分配的内存可以容纳的元素个数；
    \item \verb|vec.empty()| 返回顺序表是否为空；
    \item \verb|vec.resize(n)| 改变顺序表的大小为 \verb|n|，如果 \verb|n| 大于当前大小，则新元素默认初始化为 0；如果 \verb|n| 小于当前大小，则多余的元素被删除；
    \item \verb|vec.reserve(n)| 改变顺序表的容量为 \verb|n|，如果 \verb|n| 大于当前容量，则分配新的内存；如果 \verb|n| 小于当前容量，则不做任何操作；
    \item \verb|vec.clear()| 清空顺序表的所有元素，但不改变容量；
    \item \verb|vec.shrink_to_fit()| 将顺序表的容量缩小到当前大小，以节省内存。一般不使用这个，而是在表的内存利用率小于一半时，将表的容量缩小为当前大小的一半，保证内存利用率略大于一半但又不会频繁扩容。
\end{itemize}

\paragraph{检查} 我们经常会需要检查某容器内的某个或全部元素是否具有某种性质。以判断全部数值都为奇数为例，以前的写法可能是这样的：
\begin{minted}{cpp}
bool is_all_even(const std::vector<int>& vec) {
    for (const auto& x : vec) {
        if (x % 2 != 0) return false;
    }
    return true;
}
\end{minted}
在有了STL算法之后，我们可以用更简洁的写法：
\begin{minted}{cpp}
bool is_all_even(const std::vector<int>& vec) {
    return std::all_of(vec.begin(), vec.end(), [](int x) { return x % 2 == 0; });
    // return std::ranges::all_of(vec, [](int x) { return x % 2 == 0; }); // C++20 约束算法
}
\end{minted}
这里的 \verb|std::all_of| 是一个算法，它接受三个参数：一个迭代器范围（起始和结束），以及一个谓词（一个返回布尔值的函数或lambda表达式）。它会对范围内的每个元素应用谓词，如果所有元素都满足谓词，则返回 \verb|true|，否则返回 \verb|false|。很多STL算法都符合这个特征：这个迭代器范围是左闭又开区间，即包含起始迭代器指向的元素，但不包含结束迭代器指向的元素。这个迭代器范围也可以不是全部元素，而是一个子范围。比如说，如果我们只想在 \verb|vec[2]| 到 \verb|vec[5]| 之间检查是否符合这个标准，可以这样写：
\begin{minted}{cpp}
bool is_all_even(const std::vector<int>& vec, int x) {
    return std::all_of(vec.begin() + 2, vec.begin() + 6, [](int x) { return x % 2 == 0; });
}
\end{minted}

而约束算法就不需要指定迭代器范围了，它会自动使用容器的全部元素作为范围。如果想要指定范围，那需要结合 \verb|std::ranges::subrange| 或者 \verb|std::views::drop|、\verb|std::views::take| 等视图来实现：
\begin{minted}{cpp}
bool is_all_even(const std::vector<int>& vec, int x) {
    auto subrange = std::ranges::subrange(vec.begin() + 2, vec.begin() + 6);
    return std::ranges::all_of(subrange, [](int x) { return x % 2 == 0; }) != subrange.end();
}
bool is_all_even2(const std::vector<int>& vec, int x) {
    return std::ranges::all_of(
        vec | std::views::drop(2) 
            | std::views::take(4), 
        [](int x) { return x % 2 == 0; });
} // drop(2) 表示跳过前两个元素，take(4) 表示取接下来的四个元素
\end{minted}

之后的算法，我们不再提及以上取子区间或子视图的写法，都按照整个容器来处理，读者可以自行指定范围。此后也不会再提及约束算法的写法，读者可以自行模仿上述代码的实现。

与之类似的还有 \verb|std::any_of|（至少一个元素满足谓词）、\verb|std::none_of|（没有元素满足谓词）等。

\paragraph{复制和移除操作} 将一个容器复制到另一个容器中，可以使用 \verb|std::copy| 算法：
\begin{minted}{cpp}
std::vector<int> vec1 = {1, 2, 3, 4, 5};
std::vector<int> vec2(vec1.size()); // 创建一个大小相同的容器
std::copy(vec1.begin(), vec1.end(), vec2.begin()); // 将 vec1 的元素复制到 vec2
// std::copy_n(vec1.begin(), 3, vec2.begin()); // 复制前3个元素
// std::copy_if(vec1.begin(), vec1.end(), vec2.begin(), [](int x) { return x % 2 == 0; }); // 复制所有偶数元素
\end{minted}
这一招可以用于快速打印数组，因为即使容器并不相同也可以直接复制：
\begin{minted}{cpp}
std::vector<int> vec = {1, 2, 3, 4, 5};
std::copy(vec.begin(), vec.end(), std::ostream_iterator<int>(std::cout, " ")); // 输出 1 2 3 4 5
\end{minted}
\verb|std::ostream_iterator<int>| 是一个输出流的迭代器，它会将元素输出到指定的输出流（这里是 \verb|std::cout|），并在每个元素后面添加一个指定的分隔符（这里是空格）。这样可以方便的将容器中的元素输出到控制台或文件中。

将一个容器中的元素移除，可以使用 \verb|std::remove| 算法：
\begin{minted}{cpp}
std::vector<int> vec = {1, 2, 3, 4, 5};
auto new_end = std::remove(vec.begin(), vec.end(), 3); // 移除所有等于3的元素，但不会真的移除，而是将后面的元素前移，返回新的逻辑结尾迭代器
std::println("{}", vec); // 输出{1, 2, 4, 5, 5}，3被覆盖了，但最后一个元素仍然存在。std::println 是C++23
std::erase(new_end, vec.end()); // 真正移除逻辑结尾之后的元素
std::println("{}", vec); // 输出{1, 2, 4, 5}，3被真正移除了
\end{minted}
因此，\verb|std::remove| 经常会和 \verb|std::erase| 一起使用，形成“擦除-移除惯用法”（Erase-Remove Idiom）。它的时间复杂度为 $O(n)$，其中 $n$ 是容器的大小。

此外，还可以使用 \verb|remove_if| 来移除满足某个条件的元素：
\begin{minted}{cpp}
std::vector<int> vec = {1, 2, 3, 4, 5};
std::remove_if(vec.begin(), vec.end(), [](int x) { return x % 2 == 0; }); // 移除所有偶数元素
std::erase(new_end, vec.end()); // 真正移除逻辑结尾之后的元素
std::println("{}", vec); // 输出{1, 3, 5}，所有偶数被移除了
\end{minted}

\begin{remark}{迭代器失效}
    注意！如果要移除满足某条件的元素，绝对不可以这么写：
    \begin{minted}{cpp}
    for (auto it = vec.begin(); it != vec.end(); ++it) 
        if (*it % 2 == 0) 
            vec.erase(it); // 错误！erase 会使迭代器失效
    \end{minted}
    这是因为，不论是 \verb|std::remove| 还是 \verb|vec.erase(it)|，都会使迭代器失效。迭代器失效指的是迭代器不能再正确的指向容器中的元素，这可能会导致未定义行为。以下操作都可能会导致迭代器失效：
    \begin{itemize}
        \item 插入元素，尤其是在中间或开头插入元素；
        \item 删除元素，尤其是在中间或开头删除元素；
        \item 容器扩容或缩容；
        \item 清空元素。
    \end{itemize}
    因此，在使用迭代器时，一定要注意它们的有效性，避免在迭代器失效后继续使用它们。
\end{remark}

\begin{remark}{bool数组}
    \verb|std::vector<bool>| 的实现非常特殊。它实际上是一个 \verb|bitset|，每个元素只占用一位，而不是一个字节。这种实现方式可以节省内存，但也带来了一些问题：
    \begin{itemize}
        \item 不能直接获取元素的引用，因为每个元素只占用一位，无法返回一个可以修改的引用，单个位也没有地址
        \item 不能直接使用迭代器，因为迭代器需要返回一个可以修改的引用，而 \verb|std::vector<bool>| 的迭代器返回的是一个代理对象，不能直接修改元素；
        \item 线程不安全，位压缩导致读写冲突可能性极高。
        \item 排序、查找等算法能用，但是缓慢。
    \end{itemize}
    因此，在空间允许的情况下，可以使用其他的数组来代替这个实现。
\end{remark}


\paragraph{查找操作} 在查找某容器中是否存在某元素时，朴素的写法是：
\begin{minted}{cpp}
bool contains(const std::vector<int>& vec, int x) {
    for (const auto& y : vec) {
        if (y == x) return true;
    }
    return false;
}
\end{minted}
为了简化，可以使用 \verb|std::find| 算法：
\begin{minted}{cpp}
bool contains(const std::vector<int>& vec, int x) {
    return std::find(vec.begin(), vec.end(), x) != vec.end();
    // return std::ranges::find(vec, x) != vec.end(); // C++20 约束算法
}
\end{minted}
这些算法都会寻找第一个满足特定条件的元素。

如果要寻找最后一个，使用 \verb|std::find_last| 。如果要查找所有（并返回一个新的序列），使用 \verb|copy_if|：
\begin{minted}{cpp}
const std::vector<int> vec = {1, 2, 3, 4, 5};
std::vector<int> result;
std::copy_if(vec.begin(), vec.end(), std::back_inserter(result), [](int x) { return x % 2 == 0; });
// result 现在包含 {2, 4}
\end{minted}
上文中的 \verb|std::back_inserter(result)| 是一个迭代器适配器，它会将元素插入到 \verb|result| 的末尾。它的作用是将算法的输出直接写入到一个容器中，而不需要手动管理容器的大小和位置。

如果要查找的是一个序列而不是单个元素，可以使用 \verb|std::search|：
\begin{minted}{cpp}
const std::vector<int> vec = {1, 2, 3, 4, 5};
const std::vector<int> subseq = {3, 4}; 
bool found = (std::search(vec.begin(), vec.end(), subseq.begin(), subseq.end()) != vec.end());
\end{minted}
当这个序列是同一元素重复$n$次时，可以使用 \verb|std::search_n|：
\begin{minted}{cpp}
const std::vector<int> vec = {1, 2, 3, 3, 3, 4, 5};
int n = 3;
bool found = (std::search_n(vec.begin(), vec.end(), n, 3) != vec.end());
\end{minted}
当$n=2$时，可以使用 \verb|std::adjacent_find|，它也可以用来查找相邻的两个元素是否满足某种条件：
\begin{minted}{cpp}
const std::vector<int> vec = {1, 2, 3, 3, 4, 5};
bool found = (std::adjacent_find(vec.begin(), vec.end(), 3) != vec.end());
\end{minted}

此外，还可以查找一些特殊数值，比如最大值 \verb|std::max_element|、最小值 \verb|std::min_element|、最大最小值 \verb|std::minmax_element| 、第$n$大的数值 \verb|std::nth_element| 等。
\begin{minted}{cpp}
const std::vector<int> vec = {1, 2, 3, 4, 5};
int max = *std::max_element(vec.begin(), vec.end()); // max = 5
int min = *std::min_element(vec.begin(), vec.end()); // min = 1
auto [min_it, max_it] = std::minmax_element(vec.begin(), vec.end()); // min_it 和 max_it 分别是指向最小值和最大值的迭代器

std::nth_element(vec.begin(), vec.begin() + 1, vec.end()); // 找第(1+1)大的元素，即第二大的元素
int second_largest = vec.at(1); // second_largest = 4
\end{minted}
上述代码中的 \verb|std::nth_element| 算法会重新排列容器中的元素，使得第$n$大的元素位于指定位置，并且该位置左边的元素都不大于它，右边的元素都不小于它。注意，这个算法并不会完全排序容器，只是保证了第$n$大的元素的位置。

上述代码中的 \verb|auto [min_it, max_it]| 是“结构化绑定”，它允许我们将 \verb|std::pair| 解包为两个独立的变量。这个结构化绑定也可以用于 \verb|std::tuple|。

特别的，如果数组是有序的，可以使用二分查找算法 \verb|std::binary_search|、\verb|std::lower_bound|、\verb|std::upper_bound| 等，它们的时间复杂度为 $O(\log n)$，比线性查找的 $O(n)$ 更高效。

\paragraph{排序操作} 

排序操作很简单，分两种：不改变相同元素相对顺序的排序 \verb|std::stable_sort| 和改变相同元素相对顺序的排序 \verb|std::sort|。它们的时间复杂度都是 $O(n \log n)$，其中 $n$ 是容器的大小。排序操作可以使用默认的比较函数（小于号），也可以自定义比较函数。
\begin{minted}{cpp}
std::vector<int> vec = {5, 2, 3, 1, 4};
std::sort(vec.begin(), vec.end()); // 升序排序
std::sort(vec.begin(), vec.end(), std::greater<int>()); // 降序排序
std::stable_sort(vec.begin(), vec.end(), [](int a, int b) { return a % 2 < b % 2; }); // 按奇偶性排序，
// 其中奇数在前，偶数在后，相同奇偶性的元素保持相对顺序不变
\end{minted}

检查一个数组是否按升序排列，可以使用 \verb|std::is_sorted| 和 \verb|std::is_sorted_until|。
\begin{minted}{cpp}
std::vector<int> vec = {1, 2, 3, 4, 5};
bool sorted = std::is_sorted(vec.begin(), vec.end()); // sorted = true
std::vector<int> vec2 = {1, 3, 2, 4, 5};
bool sorted2 = std::is_sorted(vec2.begin(), vec2.end()); // sorted2 = false
auto it = std::is_sorted_until(vec2.begin(), vec2.end()); // it 指向第一个不满足有序的元素，即 2
std::vector<int> vec3 = {5, 4, 3, 2, 1};
bool sorted3 = std::is_sorted(vec3.begin(), vec3.end(), std::greater<int>()); // sorted3 = true，按降序排列
\end{minted}

\paragraph{数值操作和计算} STL还提供了一些数值操作的方法。用的较多的有四个：生成序列、计数、累加、“求导”（差分）。

生成序列，顾名思义，就是生成一个数值序列。可以使用 \verb|std::iota| 算法，它会将一个范围内的元素赋值为从指定起始值开始的连续整数。
\begin{minted}{cpp}
std::vector<int> vec(5);
std::iota(vec.begin(), vec.end(), 1); // vec 现在包含 {1, 2, 3, 4, 5}
\end{minted}

计数也是很简单的，可以使用 \verb|std::count| 算法，它会计算一个范围内等于指定值的元素个数；也可以使用 \verb|std::count_if| 算法，它会计算一个范围内满足指定条件的元素个数。
\begin{minted}{cpp}
std::vector<int> vec = {1, 2, 3, 4, 5, 2, 3, 4};
int count_2 = std::count(vec.begin(), vec.end(), 2); // count_2 = 2
int count_even = std::count_if(vec.begin(), vec.end(), [](int x) { return x % 2 == 0; }); // count_even = 4
\end{minted}

累加，包含“部分和”和“总和”。可以使用 \verb|std::partial_sum| 算法，它会计算一个范围内的部分和，并将结果存储在另一个范围中；也可以使用 \verb|std::accumulate| 算法，它会计算一个范围内的总和，并返回结果。
\begin{minted}{cpp}
std::vector<int> vec = {1, 2, 3, 4, 5};
std::vector<int> partial_sum(vec.size());
std::partial_sum(vec.begin(), vec.end(), partial_sum.begin()); // partial_sum 现在包含 {1, 3, 6, 10, 15}
int total_sum = std::accumulate(vec.begin(), vec.end(), 0); // total_sum = 15
\end{minted}

求导（差分），用于计算一个序列的相邻元素之差。可以使用 \verb|std::adjacent_difference| 算法，它会计算一个范围内的相邻元素之差，并将结果存储在另一个范围中。
\begin{minted}{cpp}
std::vector<int> vec = {1, 2, 4, 7, 11};
std::vector<int> diff(vec.size());
std::adjacent_difference(vec.begin(), vec.end(), diff.begin()); // diff 现在包含 {1, 1, 2, 3, 4}
\end{minted}


\begin{example}{股票买卖}{}
    给定一数组，表示一只股票在不同时间的价格。你初始不持有股票，且在股市收盘时必须卖出所有股票。你至多同时持有1股股票，并可以买入/卖出至多$n$次。请设计一个算法，计算你能获得的最大利润。鉴于不同的股票管理规则，你应该设计两种算法：
    \begin{itemize}
        \item $n=1$，即只能买卖一次；
        \item $n=\infty$，即可以买卖任意次。
    \end{itemize}
    请使用STL算法，使代码简洁。

    \begin{description}
        \item[输入] 一个整数数组，表示股票在不同时间的价格。
        \item[输出] 一个整数，表示能获得的最大利润。如果没有任何交易可以获得利润，返回0。
    \end{description}
    \tcblower

    先分析“只交易一次”的情况。假设我们希望在第$i$天卖出股票，那么我们必须在第$i$天之前的某一天买入股票。为了最大化利润，我们应该在第$i$天之前选择价格最低的一天买入股票。于是，我们会得到一个新的数组：
    \begin{minted}{cpp}
    std::vector<int> prices = {7, 1, 5, 3, 6, 4}; // 假设的输入
    std::vector<int> profits(prices.size());
    for (std::size_t i = 0; i < prices.size(); ++i) {
        int min_price = *std::min_element(prices.begin(), prices.begin() + i + 1);
        profits.at(i) = prices.at(i) - min_price;
    }
    int max_profit = *std::max_element(profits.begin(), profits.end());
    \end{minted}
    现在试着优化一下这些代码。我们发现，每一次计算 \verb|std::min_element| 都重复遍历了数组。我们可以考虑一个优化：我们已知一个数组的最小值，如果在这个数组末端添加一个元素，那么添加后的数组的最小值要么是原数组的最小值，要么是新添加的元素。于是我们可以用一个变量来记录当前的最小值，以避免重复计算。同理，也可以用一个变量以类似的方式来记录当前的最大利润。于是我们可以得到一个更高效的算法：
    \begin{minted}{cpp}
    std::vector<int> prices = {7, 1, 5, 3, 6, 4}; // 假设的输入
    int min_price = std::numeric_limits<int>::max(); // 也可以写成一个极大的数，这里是取了int的上极限
    int max_profit = 0;
    for (const auto& price : prices) {
        min_price = std::min(min_price, price);
        max_profit = std::max(max_profit, price - min_price);
    }
    \end{minted}
    最后返回 \verb|max_profit| 即可。

    现在分析“可以交易任意多次”的情况。这个应当很好理解，我们可以每一天都考虑是否买入或卖出股票，只要今天的价格低于明天的价格，我们就买入股票并在明天卖出；否则，我们就不进行任何操作。于是我们可以得到一个简单的算法：
    \begin{minted}{cpp}
    std::vector<int> prices = {7, 1, 5, 3, 6, 4}; // 假设的输入
    int max_profit = 0;
    for (std::size_t i = 1; i < prices.size(); ++i) 
        if (prices.at(i) > prices.at(i - 1)) 
            max_profit += prices.at(i) - prices.at(i - 1);
    \end{minted}
    最后返回 \verb|max_profit| 即可。
\end{example}

\subsection{双端队列和栈、队列}

\paragraph{双端队列}

双端队列是顺序表的另一个实现。不同于 \verb|std::vector|，它允许在两端进行高效的插入和删除操作。双端队列的时间复杂度为：
\begin{itemize}
    \item 访问指定索引的元素所用的时间复杂度为 $O(1)$，和之前一样，因为它也是顺序表；
    \item 在双端队列的两端添加或删除一个元素所用的时间复杂度为 $O(1)$。对应的算法是 \verb|dq.push_front(x)|、\verb|dq.push_back(x)|、\verb|dq.pop_front()|、\verb|dq.pop_back()|，也有 \verb|emplace_front(x)|、\verb|emplace_back(x)|，此处它们的作用类似 \verb|std::vector| 的同名方法。
    \item 在双端队列的中间插入或删除一个元素所用的时间复杂度为 $O(n)$，和之前一样，因为它仍然需要移动后续的元素。
    \item 一般情况下，使用双端队列时，我们只会在两端进行操作，而不会在中间进行操作。访问两端的元素需要的时间复杂度是 $O(1)$，对应的算法是 \verb|dq.front()| 和 \verb|dq.back()|。
\end{itemize}

双端队列由于是顺序表，因此它也支持迭代器和STL算法。它的迭代器是随机访问迭代器，因此可以使用所有的STL算法。

其实双端队列的操作很“骚”。我们有定理：不存在一种数据结构，它既支持基于物理内存相对位置的$O(1)$随机访问，又支持在中间$O(1)$的插入或删除元素。但是，双端队列却可以在两端$O(1)$的插入或删除元素，并且支持基于物理内存相对位置的$O(1)$随机访问，这是怎么做出来的？

这是因为，双端队列并不是完好的顺序表。其核心包括两部分：
\begin{description}
    \item[缓冲区] 实际存放元素的小顺序表，每个缓冲区大小固定，比如8个\verb|int|。
    \item[中控器] 一个指针数组，每个指针指向一个缓冲区。中控器的大小是动态变化的，当缓冲区满了时，中控器会扩容；当缓冲区空了时，中控器会缩容。
\end{description}
双端队列维护两个迭代器（指针）：一个指头（第一个有效缓冲区的起始）、一个指尾（最后一个有效缓冲区的末尾）。

当我们试图基于下标访问的时候，内部做了两次计算：先算出所在缓冲区的索引（一次除法），再算出在缓冲区内的偏移量（取模），于是就可以在$O(1)$的时间内访问到元素了。当在两端插入或删除元素时，只需要移动那两个指着头和尾的指针即可；块要是满了就直接加新的块，总体时间复杂度非常低。但一旦试图在中间插入或删除，那就完蛋了：需要移动大量的元素，时间复杂度会退化到$O(n)$。

也正因此，双端队列对比顺序表的问题在于其常数较大，因为块不连续，对缓存不友好。

\paragraph{栈}

虽然双端队列足够好用，但对于大多数编程任务而言，我们更关心抽象出来的两个数据结构：栈和队列。

我们在讲函数的栈帧时就提到过了栈。这里的栈和内存上的“栈”不是一个概念，虽然它们的行为很像：栈都是后进先出的数据结构，可以想象为一个极长但极窄的管道，数据只能从一端进出。也正因此，作为数据结构的栈只能访问一端：
\begin{itemize}
    \item \verb|stack.push(x)|：将元素 \verb|x| 压入栈顶，也有类似的 \verb|emplace(x)|，它的作用是直接在栈顶构造一个 \verb|T| 对象；
    \item \verb|stack.pop()|：弹出栈顶元素；
    \item \verb|stack.top()|：访问栈顶元素，但不弹出。
\end{itemize}
这些的时间复杂度都是常数。栈在编程上非常常用，我们将以两个题目来说明。

\begin{example}{括号配对}{}
    给你一个字符串，只包含六种字符：\verb|(|、\verb|)|、\verb|[|、\verb|]|、\verb|{|、\verb|}|。请你判断所有括号是否被正确闭合。一个括号闭合的条件是：请你判断所有括号是否被正确闭合。一个括号闭合的条件是：
    \begin{itemize}
        \item 每个左括号必须有一个对应的右括号；
        \item 这对括号中所有的括号都在该括号内闭合。
    \end{itemize}
    例如 \verb|()|、\verb|()[]{}| 和 \verb|{[()]}| 都是闭合的，而 \verb|(]|、\verb|([)]| 和 \verb|{[}| 都不是闭合的。
    \begin{description}
        \item[输入] 一个字符串，只包含六种字符：\verb|(|、\verb|)|、\verb|[|、\verb|]|、\verb|{|、\verb|}|。
        \item[输出] 如果所有括号都闭合，返回 \verb|true|；否则返回 \verb|false|。
    \end{description}
    \tcblower
    看题目，观察到 \verb|{[()]}| 中， \verb|()| 是最内层的括号，它们必须先闭合，才能闭合外层的 \verb|[]|，最后才能闭合最外层的 \verb|{}|。这说明，\textbf{最后压进去的括号必须最先闭合}，这正是栈的特性。因此，我们可以使用一个栈来解决这个问题。
    
    接下来的思维就很直接了。遍历字符串中的每个字符，如果是左括号，就将其压入栈中；如果是右括号，就检查栈顶的左括号是否与之匹配，如果匹配，就弹出栈顶的左括号；如果不匹配，或者栈为空，就返回 \verb|false|。最后，如果栈为空，说明所有括号都闭合，返回 \verb|true|；否则返回 \verb|false|。
    \begin{minted}{cpp}
    bool isValid(const std::string& s) {
        std::stack<char> stk;
        for (const auto& c : s)
            if (c == '(' || c == '[' || c == '{') {
                stk.push(c);
            else {
                if (stk.empty()) return false;
                char top = stk.top();
                if ((c == ')' && top != '(') ||
                    (c == ']' && top != '[') ||
                    (c == '}' && top != '{')) {
                    return false;
                }
                stk.pop();
            }
        }
        return stk.empty();
    }
    \end{minted}
\end{example}

\begin{example}{逆波兰表达式}{}
    逆波兰表达式（Reverse Polish Notation，RPN）是一种算术表达式的表示方法，它将运算符放在操作数的后面。例如，表达式 \verb|3 + 4| 在逆波兰表示法中写作 \verb|3 4 +|，而表达式 \verb|(1 + 2) * (3 + 4)| 在逆波兰表示法中写作 \verb|1 2 + 3 4 + *|。请你计算给定的逆波兰表达式的值。有效的运算符包括 \verb|+|、\verb|-|、\verb|*| 和 \verb|/|。每个操作数可以是整数或另一个逆波兰表达式。
    \begin{description}
        \item[输入] 一个字符串数组，表示一个逆波兰表达式。保证给出的表达式合法。
        \item[输出] 一个整数，表示逆波兰表达式的值。
    \end{description}
    \tcblower
    这个题目的思维和上一题类似：当一个运算符出现时，它必然要作用于它前面的两个操作数，而这两个操作数一定是最后被压进栈的操作数。因此，\textbf{最后压进去的操作数必须最先被运算}，这正是栈的特性。
    
    所以，遍历数组中的每个元素，如果是操作数，就将其压入栈中；如果是运算符，就从栈中弹出两个操作数，进行相应的运算，然后将结果压入栈中。最后，栈中只剩下一个元素，即为逆波兰表达式的值。
    
    以下是实现代码：
    \begin{minted}{cpp}
    int evalRPN(const std::vector<std::string>& tokens) {
        std::stack<int> stk;
        for (const auto& token : tokens) {
            if (token == "+" || token == "-" || token == "*" || token == "/") {
                int b = stk.top(); stk.pop();
                int a = stk.top(); stk.pop();
                if (token == "+") stk.push(a + b);
                else if (token == "-") stk.push(a - b);
                else if (token == "*") stk.push(a * b);
                else if (token == "/") stk.push(a / b); // 注意整数除法
            } else {
                stk.push(std::stoi(token)); // 将字符串转换为整数
            }
        }
        return stk.top();
    }
    \end{minted}
\end{example}

此外，栈还用于手动递归、防止爆栈等。

\paragraph{队列}

队列是一种先进先出的数据结构，可以想象为一个极长但极窄的管道，数据只能从一端进，从另一端出。队列的操作包括：
\begin{itemize}
    \item \verb|queue.push(x)|：将元素 \verb|x| 放入队列的末尾；
    \item \verb|queue.pop()|：移除队列的头部元素；
    \item \verb|queue.front()|：访问队列的头部元素，但不移除；
    \item \verb|queue.back()|：访问队列的末尾元素，但不移除。
\end{itemize}

队列在算法中也有相当的地位。我们将以两个题目来说明。

\begin{example}{缓存}{}
    现代CPU的运算速度显著高于它访问内存的速度，所以如果频繁访问内存，程序性能会大幅下降。为了提高性能，CPU通常会使用缓存（Cache）来存储最近访问过的数据，以便在下一次访问时能够更快地获取数据。缓存的容量有限，所以当缓存满了之后，就需要淘汰一些数据。某型号CPU的缓存采用如下策略管理：当缓存满了、且要访问的数据不在缓存中时，触发缓存不命中：将最早进入缓存的数据移除，然后将新数据放入缓存。

    现在请模拟这个过程：给你一些数，表示一些内存地址；再给你一个整数，表示缓存的容量。请你计算在访问这些内存地址时，缓存不命中多少次？

    \begin{description}
        \item[输入] 2行。第一行是一个整数数组，表示一些内存地址；第二行是一个整数，表示缓存的容量。
        \item[输出] 一个整数，表示缓存不命中的次数。
    \end{description}
    \tcblower
    这个缓存的行为显然就是一个队列。所以直接把队列搬上来就行。

    \begin{minted}{cpp}
    int cacheMisses(const std::vector<int>& addresses, int capacity) {
        std::queue<int> cache;
        std::unordered_set<int> cache_set; // 用于快速判断缓存中是否存在，只需要知道这个东西查找、插入、删除的时间复杂度都是 O(1) 即可，具体讲解见后续章节
        int misses = 0;
        for (const auto& addr : addresses) {
            if (cache_set.find(addr) == cache_set.end()) { // 缓存不命中
                ++misses;
                if (cache.size() == capacity) { // 缓存满了
                    int old_addr = cache.front();
                    cache.pop();
                    cache_set.erase(old_addr);
                }
                cache.push(addr);
                cache_set.insert(addr);
            }
        }
        return misses;
    }
    \end{minted}
\end{example}

\begin{example}{约瑟夫问题}{}
    有一群人围成一圈报数，从第一个人开始报数，报到第 $m$ 个人时，这个人出列，然后从下一个人开始重新报数。这个过程一直持续下去，直到所有人都出列。请输出一个列表，表示每个人出列的顺序。
    \begin{description}
        \item[输入] 两个整数，分别表示总人数 $n$ 和报数的上限 $m$。
        \item[输出] 一个整数数组，表示每个人出列的顺序。
    \end{description}
    \tcblower
    这个问题可以用队列来模拟。我们将所有人编号为 $0$ 到 $n-1$，然后将他们依次放入队列中。每次报数时，我们将队列的头部元素移到队列的末尾，直到报到第 $m$ 个人时，将这个人出列（即从队列中移除）。重复这个过程，直到队列为空。
    
    \begin{minted}{cpp}
    std::vector<int> josephus(int n, int m) {
        std::queue<int> q;
        for (int i = 0; i < n; ++i) q.push(i);
        std::vector<int> result;
        while (!q.empty()) {
            for (int i = 1; i < m; ++i) {
                q.push(q.front());
                q.pop();
            }
            result.push_back(q.front());
            q.pop();
        }
        return result;
    }
    \end{minted}
    当然，这个题目不使用队列也是可以做的：
    \begin{minted}{cpp}
    std::vector<int> josephus(int n, int m) {
        std::vector<int> people(n);
        std::iota(people.begin(), people.end(), 0); // 初始化编号为 0 到 n-1
        std::vector<int> result;
        int index = 0;
        while (!people.empty()) {
            index = (index + m - 1) % people.size(); // 找到第 m 个人的索引
            result.push_back(people[index]);
            people.erase(people.begin() + index); // 移除这个人
        }
        return result;
    }
    \end{minted}
\end{example}

\subsection{字符串}

字符串在C++中也是顺序表，但有点区别。

字符串使用了SSO优化：当字符串小时（小于15个字符），字符串会在栈上分配内存；当字符串大时（大于等于15个字符），字符串会在堆上分配内存。

字符串的行为和 \verb|std::vector| 类似，但有一些特殊的成员函数，比如 \verb|substr| 可以切分子串、\verb|find| 可以查找子串、\verb|replace| 可以替换子串、\verb|insert| 可以插入子串、\verb|erase| 可以删除子串等。它们的时间复杂度和 \verb|std::vector| 类似，其他方法略。

\subsection{链表}

链表也是一种线性结构。它不要求元素在内存中是连续的，而是在每一个元素上维护指向下一个元素的指针（和指向前一个元素的指针），以此来保证元素之间的逻辑顺序。根据是否维护指向前一个元素的指针，链表分为单链表 \verb|std::forward_list| 和双链表 \verb|std::list|。它们的时间复杂度为：
\begin{itemize}
    \item 访问指定索引的元素所用的时间复杂度为 $O(n)$，因为链表不支持随机访问，只能从头开始遍历。对应的操作是 \verb|lst.at(i)|。
    \item 在链表的指定索引插入、删除一个元素所用的时间复杂度是$O(n)$。插入删除元素这一项本身只需要更改两个指针的指向，时间复杂度是 $O(1)$，但是由于链表不支持随机访问，所以需要先遍历到指定索引的位置，时间复杂度是 $O(n)$。对双链表，操作是 \verb|lst.insert(it, x)|、\verb|lst.erase(it)|，其中 \verb|it| 是一个迭代器，指向要插入或删除的位置。对单链表，则只能移除某元素之后的元素，操作是 \verb|lst.insert_after(it, x)|、\verb|lst.erase_after(it)|；在单链表头部插入元素的操作是 \verb|lst.push_front(x)|，在单链表头部删除元素的操作是 \verb|lst.pop_front()|。
\end{itemize}

由于链表是线性访问迭代器，所以不支持随机访问迭代器，包括所有的排序算法和基于排序的算法（例如 \verb|std::nth_element|），以及二分查找的算法。但它们有自己的优化算法，是他们的成员函数。例如 \verb|std::list::sort|、\verb|std::list::merge|、\verb|std::list::unique| 等。使用例：
\begin{minted}{cpp}
std::list<int> lst = {5, 2, 3, 1, 4};
lst.sort(); // 升序排序
std::list<int> lst2 = {3, 4, 5, 6};
lst.merge(lst2); // 合并两个有序链表，lst 现在包含 {1, 2, 3, 4, 5, 6}，lst2 现在为空
lst.unique(); // 移除相邻的重复元素，lst 现在包含 {1, 2, 3, 4, 5, 6}
lst.sort(std::greater<int>()); // 降序排序，lst 现在包含 {6, 5, 4, 3, 2, 1}
\end{minted}

由于单链表的迭代器是单向的，所以部分公司面试时会要求考生手搓单链表的排序、反转等。
\begin{example}{单链表冒泡}{}
小王参加了一家公司的面试。这家不讲道理的公司要求小王对单链表进行\textbf{冒泡}排序，且\textbf{必须交换节点}，不能交换节点的值。小王虽然知道单链表排序的标准做法实际上是归并排序，但为了面试，他还是决定硬着头皮写冒泡排序。但他最终没能写出来，请你帮帮他吧！
\begin{description}
    \item[输入] 指向单链表的头节点的指针。
    \item[输出] 指向排序后单链表的头节点的指针。
\end{description}
这个单链表的节点是这样的：
\begin{minted}{cpp}
struct ListNode {
    int val;
    std::unique_ptr<ListNode> next;
    ListNode(int x) : val(x), next(nullptr) {}
};
\end{minted}
\tcblower
    由于单链表的迭代器是单向的，所以我们只能从头节点开始遍历。我们可以使用两个指针，分别指向当前节点和下一个节点。如果当前节点的值大于下一个节点的值，就交换它们的位置。
    由于我们需要交换节点，所以我们还需要一个指向当前节点的前一个节点的指针。我们可以使用一个虚拟头节点来简化边界情况的处理。虚拟头节点的值可以是任意值，它的下一个节点指向原链表的头节点。这样，我们就可以统一处理所有节点的交换，而不需要考虑头节点的特殊情况。
    \begin{minted}{cpp}
std::unique_ptr<ListNode> bubbleSortList(std::unique_ptr<ListNode> head) {
    if (!head || !head->next) return head; // 空链表或只有一个节点的链表不需要排序
    auto dummy = std::make_unique<ListNode>(0); // 创建虚拟头节点
    dummy->next = std::move(head); // 将原链表接到虚拟头节点后面
    bool swapped;
    do {
        swapped = false;
        auto prev = dummy.get(); // 前一个节点，初始为虚拟头节点
        auto curr = prev->next.get(); // 当前节点，初始为原链表的头节点
        while (curr && curr->next) { // 遍历链表
            if (curr->val > curr->next->val) { // 如果当前节点的值大于下一个节点的值，就交换它们
                auto next = std::move(curr->next); // 保存下一个节点
                curr->next = std::move(next->next); // 将当前节点的下一个节点指向下一个节点的下一个节点
                next->next = std::move(prev->next); // 将下一个节点的下一个节点指向当前节点
                prev->next = std::move(next); // 将前一个节点的下一个节点指向下一个节点
                swapped = true; // 标记已经交换过节点
            } else {
                curr = curr->next.get(); // 否则，移动到下一个节点
            }
            prev = prev->next.get(); // 前一个节点也要移动到下一个节点
        }
    } while (swapped); // 如果有交换过节点，就继续遍历链表
    return std::move(dummy->next); // 返回排序后的链表的头节点
}
    \end{minted}

    当然，我们实际上应该知道，单链表的排序最优解是归并排序，时间复杂度为 $O(n \log n)$，也好写得多，不需要考虑单链表的前驱。假如这个题目要求的是双链表，那也简单得多了！

    单链表的归并实际上应该这样：
    \begin{minted}{cpp}
std::unique_ptr<ListNode> mergeSortList(std::unique_ptr<ListNode> head) {
    if (!head || !head->next) return head; // 空链表或只有一个节点的链表不需要排序
    // 使用快慢指针找到链表的中点
    auto slow = head.get();
    auto fast = head.get();
    std::unique_ptr<ListNode>* prev = nullptr; // 用于断开链表
    while (fast && fast->next) {
        prev = &slow->next; // 保存慢指针的前一个节点
        slow = slow->next.get();
        fast = fast->next->next.get();
    }
    // 断开链表
    std::unique_ptr<ListNode> mid = std::move(*prev);
    *prev = nullptr; // 断开链表
    // 递归排序左右两部分
    auto left = mergeSortList(std::move(head));
    auto right = mergeSortList(std::move(mid));
    // 合并两个有序链表
    auto dummy = std::make_unique<ListNode>(0);
    auto curr = dummy.get();
    while (left && right) {
        if (left->val < right->val) {
            curr->next = std::move(left);
            curr = curr->next.get();
            left = std::move(curr->next);
        } else {
            curr->next = std::move(right);
            curr = curr->next.get();
            right = std::move(curr->next);
        }
    }
    curr->next = std::move(left ? left : right);
    return std::move(dummy->next);
}
    \end{minted}

    双链表的冒泡排序应该这样，不需要麻烦的维护一个额外的前驱指针：
    \begin{minted}{cpp}
struct ListNode {
    int val;
    std::unique_ptr<ListNode> next; 
    ListNode* prev; // 前驱
    ListNode(int x) : val(x), next(nullptr), prev(nullptr) {}
};

std::unique_ptr<ListNode> bubbleSortDoublyList(std::unique_ptr<ListNode> head) {
    if (!head || !head->next) return head; // 空链表或只有一个节点的链表不需要排序
    bool swapped;
    do {
        swapped = false;
        auto curr = head.get(); // 当前节点，初始为原链表的头节点
        while (curr && curr->next) { // 遍历链表
            if (curr->val > curr->next->val) { // 如果当前节点的值大于下一个节点的值，就交换它们
                auto next = std::move(curr->next); // 保存下一个节点
                curr->next = std::move(next->next); // 将当前节点的下一个节点指向下一个节点的下一个节点
                if (curr->next) curr->next->prev = curr; // 更新下一个节点的前驱
                next->next = std::move(curr); // 将下一个节点的下一个节点指向当前节点
                next->prev = curr->prev; // 更新下一个节点的前驱
                if (curr->prev) curr->prev->next = std::move(next); // 更新前驱的下一个节点
                else head = std::move(next); // 如果当前节点是头节点，更新头节点
                curr->prev = next.get(); // 更新当前节点的前驱
                swapped = true; // 标记已经交换过节点
            } else {
                curr = curr->next.get(); // 否则，移动到下一个节点
            }
        }
    } while (swapped); // 如果有交换过节点，就继续遍历链表
    return std::move(head); // 返回排序后的链表的头节点
}
    \end{minted}
\end{example}

\begin{example}{单链表反转}{}
    给你一个单链表的头节点，请你反转这个链表，并返回反转后的头节点。
    \begin{description}
        \item[输入] 指向单链表的头节点的指针。
        \item[输出] 指向反转后单链表的头节点的指针。
    \end{description}
    这个单链表的节点是这样的：
    \begin{minted}{cpp}
struct ListNode {
    int val;
    std::unique_ptr<ListNode> next;
    ListNode(int x) : val(x), next(nullptr) {}
};
    \end{minted}
    \tcblower
    非常容易想到的一个算法是：遍历一次链表，把链表的每个节点都压入栈中，然后再从栈中弹出节点，重新连接它们。这样就可以得到一个反转后的链表。时间复杂度为 $O(n)$，空间复杂度为 $O(n)$，时间复杂度已经是最优的了。但有没有“原地”反转的可能性呢？

    当然有。我们可以使用三个指针，分别指向当前节点、前一个节点和下一个节点。遍历链表时，将当前节点的下一个节点指向前一个节点，然后将前一个节点和当前节点都向后移动一位。这样就可以在遍历链表的同时反转链表。时间复杂度为 $O(n)$，空间复杂度为 $O(1)$，是最优解。
    \begin{minted}{cpp}
std::unique_ptr<ListNode> reverseList(std::unique_ptr<ListNode> head) {
    std::unique_ptr<ListNode> prev = nullptr; // 前一个节点，初始为 nullptr
    while (head) { // 遍历链表
        auto next = std::move(head->next); // 保存下一个节点
        head->next = std::move(prev); // 将当前节点的下一个节点指向前一个节点
        prev = std::move(head); // 将前一个节点移动到当前节点
        head = std::move(next); // 将当前节点移动到下一个节点
    }
    return std::move(prev); // 返回反转后的链表的头节点
}
    \end{minted}
\end{example}
\subsection{红黑树}

顺序表支持随机访问，但在中间插入或删除元素的时间复杂度为 $O(n)$；链表支持在中间插入或删除元素的时间复杂度为 $O(1)$，但不支持随机访问。我们希望有一种数据结构，它既支持随机访问，又支持在中间快速的插入或删除元素。

能不能把访问和插入/删除的时间复杂度都降低到$O(1)$呢？很遗憾，不能。实际上有这个定理，“不存在这样的数据结构”是这个定理的推论。
\begin{theorem}{Fredman-Saks 下界}{}
    任何数据结构维护一个动态序列，其按物理位置查找和插入/删除的时间复杂度的和至少为$\Omega(\log n / \log \log n)$，其中$n$是序列的长度。
\end{theorem}

\begin{corollary}{}{}
    不存在一种数据结构，它既支持基于物理内存位置的$O(1)$随机访问，又支持在中间$O(1)$的插入或删除元素。
\end{corollary}

所以我们只能退而求其次，利用定理的结论，让查找、插入/删除的时间复杂度的和至少为$\Omega(\log n / \log \log n)$。一般的实现是$O(\log n)$的，例如红黑树。

为了解释清楚红黑树，要先解释清楚二叉搜索树。

\paragraph{二叉搜索树} “树”是一种数据结构，它由节点组成。每一个节点都包含指向其子节点的指针，一个节点可以有任意多个子节点，也可以没有子节点；一个节点至多有一个父节点。没有父节点的节点被称为“根”节点，没有子节点的节点被称为“叶子”节点。一棵树有且仅有一个根节点，但可以有任意多个叶子节点。

插入，显然是常数的（因为只需要动一个指针），但查找和删除看起来比线性数据结构还要复杂啊！

我们不妨思考一下在\textbf{有序的}顺序表中怎么去查找一个元素。最快的方法自然是二分查找：先和中间值比较，如果小于中间值，就在左半部分继续查找；如果大于中间值，就在右半部分继续查找。每次查找都将查找范围缩小一半，直到找到目标元素或者查找范围为空。

那么，模仿这个过程，保证这个树的一个节点至多有2个子节点，左子节点比他小，右子节点比他大，那么查找的时候就可以这样：先和中间值比较，如果小于中间值，就在左子树继续查找；如果大于中间值，就在右子树继续查找。每次查找都将查找范围缩小一半，直到找到目标元素或者查找范围为空。而向树中插入一个元素时，也可以按照同样的规则找到合适的位置插入，且新插入的元素一定是叶子节点。

最后的问题自然是如何删除元素了。当删除的是叶子时，直接删除即可；删除树中间的元素时，我们可以将这个节点的值替换为它的左子树中最大的值或者右子树中最小的值，然后删除那个叶子节点。这样就可以保证树的结构不变。

这个东西，就叫\textbf{二叉搜索树}（Binary Search Tree，BST），如图~\cref{fig:bst} 所示。
\begin{minted}{cpp}
struct TreeNode {
    int val;
    std::unique_ptr<TreeNode> left;
    std::unique_ptr<TreeNode> right;
    TreeNode(int x) : val(x), left(nullptr), right(nullptr) {}
};

class BinarySearchTree {
    std::unique_ptr<TreeNode> root;
public:
    BinarySearchTree() : root(nullptr) {}
    BinarySearchTree(const std::vector<int>& values) : root(nullptr) {
        for (const auto& value : values) insert(value);
    }
    void insert(int value) {
        if (!root) {
            root = std::make_unique<TreeNode>(value);
            return;
        }
        TreeNode* current = root.get();
        while (true) {
            if (value < current->val) {
                if (!current->left) {
                    current->left = std::make_unique<TreeNode>(value);
                    return;
                }
                current = current->left.get();
            } else {
                if (!current->right) {
                    current->right = std::make_unique<TreeNode>(value);
                    return;
                }
                current = current->right.get();
            }
        }
    }
    bool search(int value) const {
        TreeNode* current = root.get();
        while (current) {
            if (value == current->val) return true;
            current = (value < current->val) ? current->left.get() : current->right.get();
        }
        return false;
    }
    void remove(int value) {
        root = removeNode(std::move(root), value);
    }
private:
    std::unique_ptr<TreeNode> removeNode(std::unique_ptr<TreeNode> node, int value) {
        if (!node) return nullptr;
        if (value < node->val) {
            node->left = removeNode(std::move(node->left), value);
        } else if (value > node->val) {
            node->right = removeNode(std::move(node->right), value);
        } else {
            if (!node->left) return std::move(node->right);
            if (!node->right) return std::move(node->left);
            TreeNode* minNode = findMin(node->right.get());
            node->val = minNode->val;
            node->right = removeNode(std::move(node->right), minNode->val);
        }
        return node;
    }
    TreeNode* findMin(TreeNode* node) const {
        while (node && node->left) node = node->left.get();
        return node;
    }
};  
\end{minted}

但是，朴素的二叉搜索树有一个问题：如果插入的元素本来就高度有序，那么这东西最终很快就会退化成链表，这样查找的时间复杂度就会从 $O(log n)$ 退化到 $O(n)$，如图~\cref{fig:bad-bst} 所示。

\begin{figure}
    \begin{subfigure}{0.45\textwidth}
        \centering
        \includestandalone{images/binary-search-tree}
        \caption{二叉搜索树，插入的顺序是537246}
        \label{fig:bst}
    \end{subfigure}
    \hfill
    \begin{subfigure}{0.45\textwidth}
        \centering
        \includestandalone{images/bad-bst}
        \caption{几乎退化成链表的二叉搜索树，插入的顺序是654378}
        \label{fig:bad-bst}
    \end{subfigure}
\end{figure}

很容易看出来到底发生了什么。当我们按照~\cref{fig:bad-bst} 中的顺序插入元素时，根据朴素的二叉树规则，新的节点总是加在左（右）边，导致树的高度变得很高，但每一层并没有“填满”；而好的二叉树应该是~\cref{fig:bst} 中的样子，树的高度尽量低，每一层尽量填满。于是我们就需要一种机制来保证二叉搜索树的高度尽量低，这就是红黑树。

\paragraph{红黑树}

红黑树是一种二叉搜索树。比起朴素的二叉搜索树，红黑树只干了一件事：设定一套规则，让树的高度尽量低、每一层尽量填满。红黑树的规则如下：
\begin{itemize}
    \item 每一个节点，要么是红的，要么是黑的；
    \item 根节点必须是黑的\footnote{部分书籍允许根节点是红的。“根节点必须是黑的”这条性质要求完成插入操作后若根节点为红色则将其染黑，但实际这个操作确实可以延迟到删除操作时再进行，所以不是本质的要求。}；
    \item 每一个叶子节点（NIL节点，空节点）是黑的；
    \item 如果一个节点是红的，那么它的两个子节点必须是黑的；
    \item 新插入的节点必须是红的（如果因此违反了规则，需要进行调整）；
    \item 对于每一个节点，从该节点到其所有后代叶子节点的简单路径上，必须包含相同数目的黑色节点。
\end{itemize}

\begin{figure}
    \centering
    \includestandalone{images/rbt}
    \caption{一棵红黑树}
    \label{fig:rbt}
\end{figure}

所以，在插入一个节点时，新节点的父亲可能是黑的，也可能是红的。
\begin{itemize}
    \item 如果新节点的父亲是黑的，那么插入操作完成后，红黑树仍然满足所有性质；
    \item 如果新节点的父亲是红的，那么插入操作完成后，红黑树不再满足性质4（红色节点不能有红色子节点），需要调整，此时包括：
    \begin{itemize}
        \item 父叔同红~\cref{fig:rbt-insert1-1}，父叔都染黑~\cref{fig:rbt-insert1-2}，调整结束。
        \item 父红叔黑，父子同向（即父子都是左孩子，或者父子都是右孩子）~\cref{fig:rbt-insert2-1}：父变祖父，祖父变父亲的右（左）子树~\cref{fig:rbt-insert2-2}，父染黑、祖染红~\cref{fig:rbt-insert2-3}，调整结束。
        \item 父红叔黑，父子异向（即父亲是左孩子，子节点是右孩子，或者父亲是右孩子，子节点是左孩子）~\cref{fig:rbt-insert3-1}：子变父，父变子的左（右）子树~\cref{fig:rbt-insert3-2}。现在变成了第二种情况（父子同向），继续调整。
    \end{itemize}
\end{itemize}

\begin{figure}
    \centering
    \begin{subfigure}{0.3\textwidth}
        \centering
        \includestandalone{images/rbt-modify1-1}
        \caption{}
        \label{fig:rbt-insert1-1}
    \end{subfigure}
    % \hfill
    \begin{subfigure}{0.3\textwidth}
        \centering
        \includestandalone{images/rbt-modify1-2}
        \caption{}
        \label{fig:rbt-insert1-2}
    \end{subfigure}
    \caption{红黑树插入节点的调整过程1：父叔同红}
\end{figure}
\begin{figure}
    \begin{subfigure}{0.3\textwidth}
        \centering
        \includestandalone{images/rbt-modify2-1}
        \caption{}
        \label{fig:rbt-insert2-1}
    \end{subfigure}
    \hfill
    \begin{subfigure}{0.3\textwidth}
        \centering
        \includestandalone{images/rbt-modify2-2}
        \caption{}
        \label{fig:rbt-insert2-2}
    \end{subfigure}
    \hfill
    \begin{subfigure}{0.3\textwidth}
        \centering
        \includestandalone{images/rbt-modify2-3}
        \caption{}
        \label{fig:rbt-insert2-3}
    \end{subfigure}
    \caption{红黑树插入节点的调整过程2：父红叔黑，父子同向}
\end{figure}
\begin{figure}
    \centering
    \begin{subfigure}{0.3\textwidth}
        \centering
        \includestandalone{images/rbt-modify3-1}
        \caption{}
        \label{fig:rbt-insert3-1}
    \end{subfigure}
    % \hfill
    \begin{subfigure}{0.3\textwidth}
        \centering
        \includestandalone{images/rbt-modify3-2}
        \caption{}
        \label{fig:rbt-insert3-2}
    \end{subfigure}
    \caption{红黑树插入节点的调整过程3：父红叔黑，父子异向}
\end{figure}

至于删除，那麻烦就更大了。首先，红黑树是二叉搜索树，所以删除节点的方式和二叉搜索树完全一致：如果删除的是叶子节点，直接删除即可；如果删除的是中间节点，则将这个节点的值替换为它的左子树中最大的值或者右子树中最小的值，然后删除那个叶子节点。然后再根据红黑树的性质进行调整，这里不妨记真正被删除的节点的颜色$y$，以及谁来“顶替”这个被删除的节点$x$。容易理解，如果$y$是红色，那么这个操作不会破坏任何红黑树的性质。但如果$y$是黑色，这个操作会破坏性质4（根到nil的黑色节点数目要相同）。所以，有以下四条情况，请读者自行验证调整后的树会满足红黑树的性质：
\begin{itemize}
    \item 兄弟是红的。此时，令兄弟成为新的父亲，父亲成为兄弟的左（$y$在左）或右（$y$在右）子树，但$y$还是父亲的儿子。然后，交换兄弟和父亲的颜色，转为下列情况之一。
    \item 兄弟和所有侄子都是黑的。此时，兄弟变红，父亲染黑，调整结束。
    \item 兄弟是黑的，更远的侄子（$y$在左，则这里是兄弟的右节点；反之亦然）是红的。此时，兄弟变父亲，父亲变成兄弟的左（$y$在左）或右（$y$在右）子树，交换兄弟和父亲的颜色，将新叔叔（原来的远侄子）染成黑色，调整结束。
    \item 兄弟是黑的，更近的侄子是红的。此时，更近的侄子变兄弟，兄弟变成更近的侄子的右（$y$在左）或左（$y$在右）子树，新兄弟（原来的侄子）变黑，新远侄子（原来的兄弟）变红，转为第三种情况继续处理。
\end{itemize}

红黑树的代码复杂，这里就不展示了。在C++中，红黑树的实现是 \verb|std::set| 和 \verb|std::map|，它们的底层实现就是红黑树。由于红黑树的高度是$O(\log n)$，所以它们的查找、插入和删除操作的时间复杂度都是$O(\log n)$。

\paragraph{集合} \verb|std::set<T>| 是一个集合，它的底层实现是红黑树。它的特点是：
\begin{itemize}
    \item 元素是唯一的，不能重复；
    \item 元素是有序的，默认是升序排列；
    \item 查找、插入和删除操作的时间复杂度都是$O(\log n)$。
\end{itemize}

\verb|std::set| 的迭代器是顺序访问迭代器（因为是树），所以不支持排序（虽然它天然“有序”），也不支持二分查找（但其底层天然是二分查找），请使用 \verb|std::set::find|、\verb|std::set::lower_bound|、\verb|std::set::upper_bound| 等成员函数来查找元素。

\textbf{但该容器的空间复杂度和时间复杂度常数非常大}，如果是偏静态的查找（构建后只查不改）且数据量不大（1000以内），其效果不如 \verb|std::vector|、\verb|std::sort| 和 \verb|std::binary_search| 的组合，这个的时间常数非常小，在不涉及频繁插删的情况下更快。

此外，\verb|std::set| 的有序性依赖于严格弱序，判断重复用的不是 \verb|==|，而是 \verb|<|。如果 \verb|a<b| 和 \verb|b<a| 都为假，则认为 \verb|a| 和 \verb|b| 是相同的元素。也就是说，\verb|std::set| 的有序性依赖于严格弱序，而不是等价关系。另外，这个比较必须是稳定且不改变元素状态的，否则会产生未定义行为。类似的，绝对禁止使用 \verb|std::set| 的迭代器去修改元素的值，因为这会破坏红黑树的性质，导致未定义行为。C++ 标准规定：修改任何元素键值若改变排序位置，属于未定义行为，编译器不会报错，但运行结果随机。

此外，\verb|std::set| 必须使用异构查找来消灭临时对象，即在 \verb|std::set<T, std::less<>>| 中使用 \verb|std::less<>|，而不是 \verb|std::less<T>|。否则，查找时会创建一个临时对象，这个临时对象的生命周期只在查找函数中，浪费一次堆内存分配。

此外，\verb|std::set| 是线程不安全的，高并发场景不要使用该容器。

最后，节点句柄使得现在拿出某一个特定的元素很容易了：
\begin{minted}{cpp}
std::set<int> s = {1, 2, 3, 4, 5};
auto it = s.extract(5); // 提取元素 5，返回一个节点句柄
s.insert(std::move(it)); // 将节点句柄重插入到集合中
\end{minted}
这样是零拷贝用法，能避免不必要的拷贝和析构，提升性能。这个拿出来的句柄甚至可以修改，但 \verb|std::set| 不能改，\verb|std::map| 等就能改了。

但 \verb|std::set| 的巨大优势在于，除非迭代器指向一个被删除的元素，否则迭代器不会失效。也就是说，\verb|std::set| 的迭代器是稳定的。

因此，\verb|std::set| 仅在一种情况下有用：频繁插入、删除、查找，且数据量极大。否则，请尽量使用其他数据结构。

\paragraph{映射} \verb|std::map<K, V>| 是一个映射，它的底层实现也是红黑树。它的特点是：
\begin{itemize}
    \item 键是唯一的，不能重复；
    \item 键是有序的，默认是升序排列；
    \item 查找、插入和删除操作的时间复杂度都是$O(\log n)$。
\end{itemize}
在业务逻辑上，把 \verb|std::map<K, V>| 看作一个字典就行。它基本上和 \verb|std::set| 一样：红黑树底层、严格弱序比较、异构查找、线程不安全、迭代器稳定、适用范围等。

但因为它多了一个 \verb|V|，所以有了一些好处和注意事项。

首先，\verb|std::map| 的迭代器是 \verb|std::pair<const K, V>|，也就是说，迭代器指向的键是 \verb|const| 的，不能修改。因为如果修改了键，那么这个元素在红黑树中的位置就会改变，从而破坏红黑树的性质，导致未定义行为。但值可以随便改，包括迭代器等。

其次，\verb|std::map::operator[]| 的行为和 \verb|std::map::at| 不同。前者如果键不存在，会创建一个新的元素，值为默认构造的 \verb|V|；后者如果键不存在，会抛出 \verb|std::out_of_range| 异常。因此，绝对禁止用方括号去查找，必须用 \verb|std::map::find|、\verb|std::map::contains| 来查找。此外，使用方括号时的“不存在则插入”性质会要求键值对中的值默认可构造，否则不能编译，但其他算法或成员方法依然可用。若想插入或修改，使用 \verb|std::map::insert_or_assign| 更好一些。

最后，\verb|std::map::emplace| 和 \verb|std::map::try_emplace|的行为不同：前者无论插入是否成功，都会去构造；而后者只有当插入成功（键不存在）时，才会去构造。也就是说，\verb|std::map::try_emplace| 可以避免不必要的构造，从而提高性能。

\paragraph{多重集合、多重映射} 这俩玩意依然是红黑树。内存模型、迭代器稳定性、异构查找、严格弱序的比较器，线程不安全、适用范围等都和 \verb|std::set|、\verb|std::map| 一样。但他们允许等价元素的重复，所以这就把许多行为彻底改写了。

首先，我们使用多重集合、多重映射的时候，往往查找时要的就不是“第一个”了，而是“全部”。那么这时候用 \verb|find| 就不行，必须使用 \verb|equal_range|，它返回一个区间，包含所有等价元素的迭代器。然后我们就可以遍历这个区间，处理所有等价元素。
\begin{minted}{cpp}
std::multiset<int> ms = {1, 2, 2, 3, 4};
auto [beg, end] = ms.equal_range(2); // 返回一个区间，包含所有等价元素的迭代器
for (auto it = beg; it != end; ++it) {
    std::cout << *it << " "; // 输出 2 2
}
\end{minted}
但如果只想知道“有没有”，不想知道“有哪些”，那就用 \verb|find| 也行。因为 \verb|find| 返回的是第一个等价元素的迭代器，如果返回的是 \verb|end|，说明没有；否则，说明有。不能 \verb|count|，它会大规模遍历、偷吃性能，只有在你想知道“有多少个”等价元素时才用 \verb|count|。

其次，这俩东西的插入操作是\textbf{永远成功}的：重复键一定允许插入，不会存在单集合、单映射的拒绝插入场景。所以，在实际业务逻辑中，我们不能通过返回值来判断“这是否首次出现”，必须使用 \verb|count| 或者 \verb|equal_range| 来判断“有多少个”等价元素。

再次，这东西不能用方括号和 \verb|at| 来查找，因为他们的键不是唯一的，不能用一个键去唯一确定一个值。要插入，必须 \verb|insert|；要修改，必须先 \verb|equal_range|，找到具体的迭代器再改（当然只能改多重映射的值）。

最后，\verb|erase| 也不一样：\verb|erase(key)| 会删除所有等价元素；\verb|erase(iterator)| 会删除指定迭代器指向的元素。

\subsection{哈希表}

再次回顾一下平衡二叉搜索树。它能做到插入删除和检索的速度都是$O(\log n)$。我们也知道，不存在一种数据结构，它既支持基于物理内存相对位置的$O(1)$随机访问，又支持在中间$O(1)$的插入或删除元素。

那么，能不能“不基于物理内存位置”来进行$O(1)$随机访问呢？答案是肯定的，这就是哈希表。哈希表的底层是一个普通的顺序表，但它上面有一个哈希函数：我们将一个输入通过哈希函数能映射到一个整数，这个整数对数组长度取模，就能得到一个数组下标，形成变相的随机访问。如果不同的键算出相同的下标（哈希冲突），可以用其他手段来解决，比如拉链法：把这个底层顺序表存的不再是数值，而是一个红黑树，形成一个红黑树森林（这里一般叫“桶”，也可以用其他数据结构来实现桶）。如果发生冲突，就把新元素挂在对应的红黑树末尾。这样，哈希表的查找、插入和删除操作的时间复杂度都是$O(1)$，但最坏情况下（所有的键全都撞在同一个位置），就可能退化到$O(\log n)$（红黑树森林）或$O(n)$（朴素拉链），所以哈希表的性能极度依赖于哈希函数的质量。

所以实际上可以看出这是一个投机取巧的手段：顺序表的随机访问是“给我第$i$个元素”，而哈希表的“随机访问”是“给我Key为‘张三’的元素”：这么一来，哈希表里，元素根本就没有“中间”这个概念！

因此所有的插入、删除、查找，只要不发生哈希冲突，就都是$O(1)$的时间复杂度：不管输入什么，先进哈希函数滚一圈变成整数，然后利用顺序表的随机访问性质就都能在常数时间内完成。

但这导致了一个致命的问题：哈希表里的元素\textbf{不可能有序}，他也不可能提供“按物理位置插入”的功能。所以哈希表的迭代器是顺序访问迭代器，不能排序。

\paragraph{无序集合和无序映射} 在C++中，\verb|std::unordered_set<T>| 和 \verb|std::unordered_map<K, V>| 就是哈希表的实现。它们的特点是：
\begin{itemize}
    \item 元素是唯一的，不能重复；
    \item 元素是无序的，不能排序；
    \item 查找、插入和删除操作的时间复杂度都是$O(1)$，但最坏情况下可能退化到$O(n)$。
\end{itemize}

无序集合和无序映射的时间复杂度常数项非常大，所以数据不大（低于100）时，可以选择顺序表来代替。

C++20之前的 \verb|std::hash<std::string>| 是一个非常糟糕的哈希函数，容易导致哈希冲突，甚至被恶意利用进行拒绝服务攻击（DoS）。C++20之后，\verb|std::hash<std::string>| 改进了哈希函数，避免了这个问题。重点是避免多个元素落入同一个桶：有东西的桶数量越多，哈希表的性能就越好。

哈希表是按需重新分配大小的，且哈希表的重分配代价巨大：所有元素重新分桶，所有迭代器全部失效。因此，如果能预知元素数量，务必预分配桶：
\begin{minted}{cpp}
constexpr std::size_t EXPECTED_MAX_SIZE = 1000000;
std::unordered_set<int> uset;
uset.reserve(EXPECTED_MAX_SIZE * 1.5); // 预分配桶，避免哈希表重分配
\end{minted}
这里乘1.5的主要原因是C++的哈希表在填满75\%时就会进行重分配，所以预分配时要留出一定的空间。

无序集合和无序映射的比较器要求是只实现相等和哈希即可。当你的类型不是预先定义的类型时，想把这种东西塞进无序集合和无序映射的键时，必须自己特化一个 \verb|std::hash|：
\begin{minted}{cpp}
struct MyType {
    int a;
    std::string b;
};

struct MyTypeHash {
    std::size_t operator()(const MyType& obj) const {
        std::size_t h1 = std::hash<int>{}(obj.a);
        std::size_t h2 = std::hash<std::string>{}(obj.b);
        return h1 ^ (h2 << 1); // 或者使用其他组合方式
    }
}; // 这是一个仿函数，可以用括号调用

std::unordered_set<MyType, MyTypeHash> mySet;
\end{minted}
此外，一切对无序集合和无序映射的排序都是不合理的。

最后，哈希表的迭代器也是相当稳定的：当且仅当重新分桶或删除元素时，迭代器才会失效。

哈希表的应用场景主要在于：不需要排序的超大量查找，且不介意偶发的最坏情况。如果需要有序的遍历、区间查询，请使用有序集合和有序映射。在量化交易这种不能接受偶尔的最坏情况的场景下，需使用其他的方法，如做平滑扩容，或使用非标准库的其他哈希表。

\paragraph{多重无序集合} 这个还算能用。

它和无序集合差不多，只是允许重复键，所以查找所有依然要用 \verb|equal_range|，插入永远成功，删除 \verb|erase(key)| 会删除所有等价元素，删除 \verb|erase(iterator)| 会删除指定迭代器指向的元素。

一个问题是：假如有一百万个元素，且都是相同的键，那么这东西就变成链表了，挂在同一个桶里，性能非常不好。遇到这种情况，建议使用无序映射，把值当成计数器或放进数组，这样查找Key还是$O(1)$。

这东西的迭代器失效规则和无序集合一样：当且仅当重新分桶或删除元素时，迭代器才会失效。特别注意：当使用\verb|equal_range|拿到区间并移除某些元素时，必须妥善保存下一个迭代器。
\begin{minted}{cpp}
auto range = umm.equal_range(key);
for (auto it = range.first; it != range.second; ) {
    if (cond) it = umm.erase(it); 
    else ++it;
}
\end{minted}

另外，多重有序集合保证插入顺序的稳定性，即先进集合的一定在左边，后进集合的在右边。但多重无序集合对此做不到，因为插入顺序无法保证。

多重无序集合的应用面很窄：仅当你需要极频繁地插入重复 Key，且从不遍历重复组、从不关心顺序、Key 的重复度极低时，才考虑这个。

最后，说一个事：\textbf{永远不要使用多重无序映射}。工业标准就是下面这个：
\begin{minted}{cpp}
std::unordered_map<Key, std::vector<Value>> myMap;
// or std::unordered_map<Key, std::list<Value>> myMap; 当Value不多时
// never std::unordered_multimap<Key, Value>
\end{minted}

\subsection{堆}

这个堆和之前那个“栈”差不多。之前说到栈，有两个：数据结构栈，和内存栈。这里也是一样的：数据结构堆，和内存堆。我们这里的堆自然是数据结构堆，因为不涉及 \verb|new|。

堆的逻辑形态是一棵完全二叉树：在二叉树的基础上，要求从上到下、从左到右依次填满每一层。不过因为完全二叉树的良好特性，所以我们完全可以在物理上用一个顺序表去存这个。理由：除了最后一层外，其他层都是满的，且最后一层节点从左到右紧密排布，这种特性可以让我们把数组的下标和树的层次结构对应起来：根节点的下标是0，对于任意一个节点，它的左孩子下标是 $2i+1$，右孩子下标是 $2i+2$，父节点下标是 $(i-1)/2$（向0舍入）。这样，不仅节省了内存，还缓存非常友好。

但仅仅是一个“堆”，我们并不能用他来做什么，所以要给他加上一些限制。比如“大根堆”：要求对于每一个父节点，它的不小于它的两个孩子节点。这样一来，根节点就是整个堆中最大的元素。类似的还有“小根堆”，要求每一个父节点不大于它的两个孩子节点，这样根节点就是整个堆中最小的元素。我们以大根堆来举例说明。

当我们在大根堆中插入一个新元素时，我们先把它放在顺序表的末尾，此时它在大根堆的最后一层。随后，逐步上浮：如果它大于它的父亲，那么就把它和父亲换位，如此不断上浮，直到它不大于它的父亲，或者它已经是根节点了。这个过程的时间复杂度是 $O(\log n)$。

当我们在大根堆中移除根节点时（我们只关心堆顶元素，不关心其他元素），我们则反过来：把最后一个元素提到堆顶（实际上由于底层是顺序表，是将顺序表最后元素和首元素换位，然后移除最后的元素，这样是$O(1)$的），然后逐步下沉：如果它小于它的两个孩子节点中较大的那个，那么就把它和那个孩子换位，如此不断下沉，直到它不小于它的两个孩子节点，或者它已经是叶子节点了。这个过程的时间复杂度也是 $O(\log n)$。

当我们从一个普通数组建立堆的时候，它的时间复杂度是 $O(n)$：从最后一个非叶子节点开始逐步向前下沉，直到根节点。为什么是 $O(n)$ 呢？因为叶子节点不需要下沉，倒数第二层的节点最多下沉一层，倒数第三层的节点最多下沉两层……根节点最多下沉 $\log n$ 层。于是总的时间复杂度就是：
\begin{equation}
\sum_{i=0}^{\log n} \frac{n}{2^{i+1}} \cdot i = O(n)
\end{equation}

建立堆后，有一个巨大的好处：如果是大根堆，那么根节点就是整个堆里最大的；如果是小跟堆，那么根节点就是整个堆里最小的。于是我们可以用堆来做一个排序算法：堆排序。堆排序的时间复杂度是 $O(n \log n)$，空间复杂度是 $O(1)$，是一种原地排序算法，做法就是建堆，然后把倒数第一个元素和第二个元素换位（因为最大的已经排好了），然后把堆的大小减一，重新下沉堆顶元素，直到堆的大小为1。堆排序是不稳定的排序算法。

当然，我们比起排序，更关心的是每一次都能把最大值（最小值）拿到队首。于是，优先队列诞生了。

\paragraph{优先队列}

优先队列就是堆。一般的优先队列是大根堆，队首元素是整个队列中最大的元素。优先队列的插入和删除操作的时间复杂度都是 $O(\log n)$，获取队首元素的时间复杂度是 $O(1)$。

如果要用小根堆，应该这么写：
\begin{minted}{cpp}
std::priority_queue<int, std::vector<int>, std::greater<int>> pq;
\end{minted}
这是模板类 \verb|std::priority_queue| 的一个实例化，三个模板参数分别是：元素类型、底层容器类型、比较函数对象类型。默认的底层容器是 \verb|std::vector|，默认的比较函数对象是 \verb|std::less<T>|，也就是大根堆。

为什么小根堆反而是 \verb|std::greater<int>| 呢？因为标准库大概是写反了，写的是父节点和子节点之间的比较：当父节点优先于子节点时，父节点要下沉（子节点要上浮）……这大概是一个错误，且延续了多年，此时已经成为了标准库的惯例了。那我们就只能背下来了。

优先队列是非常“有性格”的数据结构：不提供迭代器，不支持直接修改或移除其中的元素，只能通过插入和删除队首元素来操作。它的主要用途是需要频繁获取最大值或最小值的场景，比如任务调度、事件驱动模拟等。

\begin{example}{武士对决}{}
    假设现在有许多武士要角斗。每个武士都有一个名字和一个体力值，当两个武士相互角斗的时候，体力值较高的武士将会获胜，而体力值较低的武士会耗尽体力，并被淘汰。然而，角斗会消耗武士的体力值，因此每一次角斗后，胜者的体力值会减少等同于败者当前体力值的数值。如果有两个武士体力相等，则他们都会耗尽体力被淘汰。为了保证公平，武士们决定按照体力值从高到低的顺序进行角斗；如果有多个体力值最高的武士，那么他们会在这些武士中选择姓名字典序最靠前的两个武士进行角斗（例如有abc三个武士，则a的姓名字典序最靠前，b次之，c最靠后）。每次角斗后，胜者会回到队列中继续参与角斗，直到只剩下一个武士或者没有武士剩下为止。

    请你写一个程序，模拟这个过程，并输出最后剩下的武士的姓名和体力值。如果没有武士剩下，则输出\verb|No warrior left|。
    
    \begin{description}
        \item[输入] 第一行输入一个整数 $n$，表示武士的数量。接下来 $n$ 行，每行输入一个字符串和一个整数，分别表示武士的姓名和体力值。
        \item[输出] 输出最后剩下的武士的姓名和体力值，如果没有武士剩下，则输出\verb|No warrior left|。
    \end{description}
    \tcblower
    这个题目思路非常直接，看到这个过程就知道是优先队列了，纯考察知不知道有这么个队列，知道的话就能$O(n\log n)$。要不然指定超时$O(n^2\log n)$。

    先把所有武士放入优先队列中，每次取出两个武士进行角斗，比较他们的体力值，更新胜者的体力值，并将胜者重新放入优先队列中。重复这个过程，直到优先队列中只剩下一个武士或者没有武士剩下为止。

    \begin{minted}{cpp}
struct Warrior {
    std::string name;
    int strength;
    Warrior(std::string n, int s) : name(n), strength(s) {}
    bool operator<(const Warrior& other) const {
        if (strength == other.strength) {
            return name > other.name; // 体力相等时，姓名字典序
        }
        return strength < other.strength; // 体力值较高的优先
    }
};

std::priority_queue<Warrior> pq;
int n;
std::cin >> n;
for (int i = 0; i < n; ++i) {
    std::string name;
    int strength;
    std::cin >> name >> strength;
    pq.emplace(name, strength);
}
while(!pq.empty()) {
    Warrior w1 = pq.top(); pq.pop();
    if (pq.empty()) {
        std::println("{} {}" , w1.name, w1.strength);
        break;
    }
    Warrior w2 = pq.top(); pq.pop();
    if (w1.strength == w2.strength) {
        // 两个武士体力相等，双方都被淘汰
        continue;
    } else if (w1.strength > w2.strength) {
        w1.strength -= w2.strength;
        pq.push(w1);
    } else {
        w2.strength -= w1.strength;
        pq.push(w2);
    }
}
    \end{minted}
\end{example}

\subsection{位集合}

位集合是另一个“顺序表”。它接受一个编译期的整数作为位集合的长度，并在内存中为每一位分配一个二进制位（bit），而不是一个字节（byte）；其成员类型是一个代理，和\verb|std::vector<bool>|类似。位集合的大小是固定的，不能动态扩展。它的主要用途是节省内存空间，尤其是在需要存储大量布尔值的情况下。

位集合可以使用字符串或整数来初始化：
\begin{minted}{cpp}
std::bitset<8> bs("11001010"); // 初始化为11001010
std::bitset<4> bs2{0xA}; // 用整数初始化，0xA是1010
std::cout << bs.to_string() << std::endl; // 输出11001010
std::cout << bs2.to_string() << std::endl; // 输出1010
\end{minted}

这个东西支持的算法和先前有较大的不同。它的主要支持算法是位运算：与、或、非、异或、左移、右移；除此之外，比较需要注意的一点是它没有 \verb|begin()| 和 \verb|end()|，而是提供了 \verb|count()|、\verb|any()|、\verb|none()|、\verb|all()| 等方法来统计集合中为真的元素个数，或者判断集合中是否有为真的元素。要访问特定位，应使用方括号或 \verb|test()| 方法。要修改特定位，应使用 \verb|set()|、\verb|reset()|、\verb|flip()| 等方法。必须说明的是，这个特定位是从右往左数的，最右边的位是第0位，最左边的位是第n-1位，所以将一个64bit的位集合仅输出右面4位，如使用循环应当逆序，但更常见的是转成字符串再处理。

也可以直接全部打印。
\begin{minted}{cpp}
std::bitset<64> bs;
bs.set(0); // 设置第0位为1
bs.set(1); // 设置第1位为1
std::cout << bs << std::endl; 
std::cout << bs.to_string().substr(60) << std::endl; // 输出右面4位
\end{minted}

此外，它还能转成整数，\verb|to_ullong()|方法将位集合转成无符号长整数，\verb|to_ulong()|方法将位集合转成无符号整数。

\subsection{对和元组}

对、元组是不同类型元素的容器。对是一个固定大小为2的容器，元组是一个固定大小为n的容器。它们可以存储不同类型的元素，并且可以通过下标访问元素。对和元组都是值类型，它们的元素可以是任意类型，包括其他对和元组。

\begin{minted}{cpp}
std::pair<int, std::string> p(1, "hello");
auto [x, y] = p; // C++17结构化绑定
std::cout << x << " " << y << std::endl; // 输出1 hello
std::tuple<int, std::string, double> t(1, "hello", 3.14);
auto [a, b, c] = t; // C++17结构化绑定
std::cout << a << " " << b << " " << c << std::endl; // 输出1 hello 3.14
\end{minted}

这两个容器一般用于函数返回多个值，或者在STL中作为键值对使用，比如 \verb|std::map| 和 \verb|std::unordered_map| 的元素类型就是 \verb|std::pair<const Key, T>|。

要返回对的第一项和第二项，可以用 \verb|p.first| 和 \verb|p.second|，要返回元组的第n项，可以用 \verb|std::get<n>(t)|，这里的 \verb|n| 是一个编译期常量，表示要访问的元素的索引。所以以下写法是错误的：
\begin{minted}{cpp}
for (int i = 0; i < 3; ++i) {
    std::cout << std::get<i>(t) << std::endl; // 错误，i不是编译期常量
}
\end{minted}
而应该用提供的现成函数：
\begin{minted}{cpp}
std::apply([](auto&&... args) {
    ((std::cout << args << std::endl), ...);
}, t);
\end{minted}
或者干脆用结构化绑定处理，能更省事。一般元组也不建议写的太大，超过5个元素就不太好读了。实在不行，可以考虑用结构体代替元组。

\subsection{其他容器}

这里的“其他容器”主要是一些特殊用途的容器，即变体 \verb|variant|，可选值 \verb|optional|，任意类型 \verb|any|。

它们实际上都是一个“盒子”，可以存储不同类型的值，但每次只能存储一个值。它们的主要用途是解决函数返回值的问题：当函数可能返回不同类型的值时，可以使用这些容器来封装返回值，从而避免使用指针或引用。
\begin{itemize}
    \item \verb|std::optional|：表示一个可能存在也可能不存在的值。它可以用来表示函数的返回值可能为空的情况，或者表示一个变量可能没有被初始化的情况。
    \item \verb|std::variant|：表示一个可能是多种类型之一的值。它可以用来表示一个变量可能是多种类型之一的情况，或者表示一个函数的返回值可能是多种类型之一的情况。
    \item \verb|std::any|：表示一个可以存储任意类型的值。它可以用来表示一个变量可以存储任意类型的值，或者表示一个函数的返回值可以是任意类型的值。
\end{itemize}

比如说：
\begin{minted}{cpp}
std::optional<int> opt; // 默认构造，opt为空
opt = 42; // 赋值，opt现在包含一个int值42
if (opt) { // 检查opt是否包含值
    std::cout << *opt << std::endl; // 输出42
}
std::variant<int, std::string> var; // 默认构造，var为空
var = 42; // 赋值，var现在包含一个int值42
var = "Hello"; // 赋值，var现在包含一个string值"Hello"
if (std::holds_alternative<int>(var)) { // 检查var是否包含int类型的值
    std::cout << std::get<int>(var) << std::endl; // 输出42
}
std::any a; // 默认构造，a为空
a = 42; // 赋值，a现在包含一个int值42
a = std::string("Hello"); // 赋值，a现在包含一个string值"Hello"
if (a.type() == typeid(int)) { // 检查a是否包含int类型的值
    std::cout << std::any_cast<int>(a) << std::endl; // 输出42
}
\end{minted}
它们因为有类型擦除的开销，所以一定要谨慎使用，尤其是在性能敏感的场景下。

\section{视图和范围}

刚刚我们讨论的都是容器本身。要处理一个容器，可能要复制或移动，或创建新的容器，这样会有额外的开销。C++20引入了视图（View）的概念，它是一种轻量级的、不可变的、惰性求值的容器适配器。视图不会存储数据，而是提供了一种访问和操作现有容器数据的方式。

既然它是“观察数据的容器适配器”，那就天然解决了以下问题：
\begin{itemize}
    \item 告别裸指针+长度这种危险的传参组合；
    \item 告别复制或移动容器的开销，实现零拷贝的性能优化；
    \item 避免 \verb|T&| 的过度约束，让函数可以接受任意类型的容器，而不仅仅是特定类型的容器，统一接口。
\end{itemize}

必须注意，无论是什么视图，其生命周期都不会比被观察的容器长。

\subsection{一般视图}

一般的视图是 \verb|std::span|。这种视图更多的是作为函数参数存在，比如说：
\begin{minted}{cpp}
void print(std::span<const int> s) {
    std::cout << "size: " << s.size() << std::endl;
    std::ranges::copy(s, std::ostream_iterator<int>(std::cout, " "));
}

int arr[] = {1, 2, 3, 4, 5};
print(arr); // 传入数组，可以；
std::vector<int> vec = {6, 7, 8, 9, 10};
print(vec); // 传入vector，也可以；
std::array<int, N> arr2 = {11, 12, 13, 14, 15};
print(arr2); // 传入array，也可以；
\end{minted}

由于视图提供了 \verb|subspan()| 方法和相关的迭代器，我们可以从一个视图中创建一个子视图，并无缝的和STL算法集成。
\begin{minted}{cpp}
std::span<const int> s1(arr); // 创建一个视图，观察数组 arr
std::span<const int> s2 = s1.subspan(1, 3); // 创建一个子视图，观察 s1 的第 1 个元素开始的 3 个元素
std::ranges::copy(s2, std::ostream_iterator<int>(std::cout, " ")); // 输出 2 3 4
\end{minted}

总而言之，只要一个函数需要读取或修改一段连续的数据，但不关心数据来自什么容器，就都可以用 \verb|std::span| 来作为参数类型。本质上可以和 \verb|T&| 相提并论，且更灵活。只有当函数需要对原来的容器本身（而非其中的内容）进行操作时，才需要传入容器本身。

\subsection{字符串视图}

\verb|std::string_view| 可以当成 \verb|std::span| 的字符串特化版本。需要注意的是，它更像是 \verb|const std::string&|，即不可变。它可以接受C字符串，也可以接受其他的字符串类型，如 \verb|std::string| 等。

在普通视图的基础上，字符串视图提供了更多的字符串操作方法，如 \verb|substr()|、\verb|find()|、\verb|rfind()|、\verb|starts_with()|、\verb|ends_with()| 等。它们的时间复杂度和对应的 \verb|std::string| 方法相同。

\subsection{视图操作（范围库）}

C++20为视图的操作提供了一整套操作手段，它们被整合在 \verb|<ranges>| 头文件中，统称为范围库（Ranges Library）。范围库提供了许多视图适配器，可以对视图进行各种操作，如过滤、转换、连接、分割等。我们先前讨论的“约束算法”，实际上就是范围库的一部分，它们可以直接作用于容器（发生了一次隐式转换）。但范围库的视图操作更为灵活，可以对视图进行链式操作，从而实现复杂的数据处理流程。

\verb|std::ranges| 是这个范围库的大框架，而其下有很多 \verb|std::ranges::views|（也可以简写成 \verb|std::views|）的视图适配器对数据进行观察和变换。它们可以通过管道操作符来进行链式调用，从而实现对视图的流水线操作。

\begin{minted}{cpp}
std::vector<int> vec = {1, 2, 3, 4, 5, 6, 7, 8, 9, 10};
auto view = vec | std::views::filter([](int x) { return x % 2 == 0; }) // 过滤出偶数
                | std::views::transform([](int x) { return x * x; }); // 将偶数平方
std::ranges::copy(view, std::ostream_iterator<int>(std::cout, " ")); // 输出 4 16 36 64 100
\end{minted}
